\chapter{Pengurangan}\label{ch:g2:subtraction}

Ketika satuan di baris atas terlalu sedikit untuk dikurangi, kita
membuka satu puluhan --- itulah \emph{pinjaman}, cerminan dari
simpanan. Bab ini menjelaskannya dengan lidi, lalu melatih
pengurangan bersusun, dengan membilang maju sebagai jalan pintas
yang ramah.

\section{Mengapa kita meminjam}

\begin{example}[Membuka ikatan puluhan]\label{ex:g2:subtraction:unbundle}
$42 - 17$: kita ingin mengambil $7$ satuan dari $2$ satuan --- terlalu
sedikit! Maka bukalah satu dari $4$ puluhan itu menjadi sepuluh
satuan lepas: sekarang baris atas punya $3$ puluhan dan $12$ satuan.
Kurangkan: $12 - 7 = 5$ satuan, dan $3 - 1 = 2$ puluhan. Jadi
$42 - 17 = 25$.
\end{example}

\begin{omfigure}
\begin{tikzpicture}[scale=0.75]
  % satu puluhan yang terikat
  \foreach \i in {0,...,9} \draw[omExo, very thick]
    (0.2*\i,0) -- (0.2*\i,1.7);
  \draw[omExo, thick, rounded corners=3pt]
    (-0.2,0.62) rectangle (2.0,1.05);
  \node[below, font=\small] at (0.9,-0.25) {satu puluhan};
  % lalu dibuka menjadi sepuluh satuan lepas
  \draw[-Stealth, omThm, very thick] (2.6,0.85) -- (3.6,0.85);
  \foreach \i in {0,...,9} \draw[omDef, very thick]
    (4.2+0.34*\i,0) -- (4.2+0.34*\i,1.7);
  \node[below, font=\small] at (5.73,-0.25) {sepuluh satuan lepas};
  \node[right, font=\small, align=left] at (8.0,0.85)
    {bukalah satu puluhan agar cukup\\satuan untuk diambil};
\end{tikzpicture}

{\small Pinjaman adalah membuka ikatan: satu puluhan menjadi sepuluh
satuan, sehingga pengurangannya bisa berjalan.}
\end{omfigure}

\section{Cara bersusun}

\begin{method}[Pengurangan bersusun]\label{met:g2:subtraction:column}
\begin{enumerate}
  \item Tulislah bilangan yang lebih besar di atas, rata di kanan;
  \item kurangkan satuannya; jika angka di atas terlalu kecil,
        pinjamlah satu puluhan (yang menjadi sepuluh satuan di kolom
        itu);
  \item kurangkan puluhannya, lalu ratusannya;
  \item \emph{periksalah}: hasilnya ditambah bilangan yang diambil
        harus mengembalikan bilangan di baris atas.
\end{enumerate}
\end{method}

\begin{example}\label{ex:g2:subtraction:worked}
Hitunglah $503 - 168$:
\[
\begin{array}{r}
5\,0\,3 \\
-\ 1\,6\,8 \\
\hline
3\,3\,5
\end{array}
\]
Satuan: $3 - 8$ terlalu kecil; pinjamlah lewat puluhan (di sana tidak
ada, jadi pinjamlah dari ratusan dulu): baris atas menjadi $4$
ratusan, $9$ puluhan, $13$ satuan. Lalu $13 - 8 = 5$; $9 - 6 = 3$;
$4 - 1 = 3$. Periksa: $335 + 168 = 503$. \checkmark
\end{example}

\section{Membilang maju}

\begin{method}[Mengurangkan dengan membilang maju]\label{met:g2:subtraction:countup}
Untuk bilangan yang berdekatan, naiklah dari yang kecil ke yang
besar lalu jumlahkan lompatannya:
\[
83 - 78: \quad 78 \to 80 \text{ adalah } 2, \ 80 \to 83 \text{ adalah } 3,
\text{ jadi } 83 - 78 = 5 .
\]
Membilang maju lebih cepat daripada meminjam ketika kedua bilangan
berdekatan --- dan begitulah persis cara pedagang memberi kembalian.
\end{method}

\begin{omfigure}
\begin{tikzpicture}[scale=0.85]
  \draw[-Stealth] (76.6,0) -- (84.4,0);
  \foreach \i in {77,...,84} \draw (\i,-0.1) -- (\i,0.1);
  \foreach \i in {78,80,83}
    \node[below, font=\footnotesize] at (\i,-0.12) {\i};
  \draw[-Stealth, omDef, thick] (78,0.3) .. controls (79,0.9) ..
    (80,0.3);
  \node[omDef, font=\small] at (79,1.0) {$+2$};
  \draw[-Stealth, omThm, thick] (80,0.3) .. controls (81.5,1.05) ..
    (83,0.3);
  \node[omThm, font=\small] at (81.5,1.1) {$+3$};
\end{tikzpicture}

{\small $83 - 78$ dengan membilang maju: dua lompatan, $2 + 3 = 5$.}
\end{omfigure}

\begin{example}[Memberi kembalian]\label{ex:g2:subtraction:change}
Sebuah mainan berharga $34$; kamu membayar dengan $50$.
Kembaliannya: bilanglah maju, $34 \to 40$ adalah $6$, $40 \to 50$
adalah $10$, jadi kembaliannya $16$. Pengurangan $50 - 34 = 16$,
dikerjakan dengan cara pedagang.
\end{example}

\section{Latihan}

\begin{exercise}[$\star$]\label{exo:g2:subtraction:1}
Hitunglah (tanpa pinjaman): $58 - 23$; \quad $176 - 45$; \quad
$389 - 164$.
\end{exercise}

\begin{exercise}[$\star$]\label{exo:g2:subtraction:2}
Hitunglah dengan pinjaman: $52 - 17$; \quad $73 - 28$; \quad
$60 - 24$.
\end{exercise}

\begin{exercise}[$\star$]\label{exo:g2:subtraction:3}
Hitunglah secara bersusun lalu periksa satu per satu: $324 - 167$;
\quad $500 - 236$; \quad $703 - 158$.
\end{exercise}

\begin{exercise}[$\star$]\label{exo:g2:subtraction:4}
Hitunglah dengan membilang maju: $62 - 58$; \quad $91 - 85$; \quad
$100 - 96$.
\end{exercise}

\begin{exercise}[$\star$]\label{exo:g2:subtraction:5}
Hitunglah kembaliannya (bilang maju): bayar $50$ untuk buku seharga
$42$; bayar $20$ untuk mainan seharga $13$; bayar $100$ untuk tas seharga $75$.
\end{exercise}

\begin{exercise}[$\star$]\label{exo:g2:subtraction:6}
``Sebatang pohon tingginya $156$ cm. Tahun ini ia tumbuh $28$ cm.
Berapa tingginya tahun lalu?'' Tulislah pengurangan dan jawabannya.
\end{exercise}

\begin{exercise}[$\star$]\label{exo:g2:subtraction:7}
Lengkapilah dengan penjumlahan kembarannya untuk memeriksa:
$85 - 39 = 46$; benarkah? $214 - 68 = 156$; benarkah?
\end{exercise}

\begin{exercise}[$\star$]\label{exo:g2:subtraction:8}
Untuk setiap cerita, pilihlah $+$ atau $-$: ``$47$ burung, $19$
terbang pergi''; ``$47$ apel, $19$ dipetik lagi''; ``Sandi punya
$47$, Lia punya $19$: berapa lebihnya Sandi?''.
\end{exercise}

\begin{exercise}[$\star$]\label{exo:g2:subtraction:9}
Isilah yang kosong: $63 - \;?\; = 45$. (Bilanglah maju dari $45$ ke $63$.)
\end{exercise}

\begin{exercise}[$\star\star$]\label{exo:g2:subtraction:10}
Sebuah permainan: pikirkan sebuah bilangan, tambahkan $30$, lalu
kurangi $18$. Hasilnya $58$. Berapa bilangan tadi? (Kerjakan mundur!)

\end{exercise}
